\section{Discussion}\label{sec:discussion}

The calibrated model produces a sharp regime split. The same instrument trades within a few basis points of par for ninety-nine percent of hours and produces a fat-tailed cross-section of price realisations for the remaining one percent. The structural device that delivers this split is the regime-dependent dispersion of the prior on the issuer fundamental, combined with a shared microstructure block. The fixed point is endogenous and the cross-section of type thresholds is heterogeneous. At the calibrated parameters, however, the per-type conversion rates in the stress regime are nearly uniform across types (within about one percentage point of each other), and the corridor differentiation in the cost comparison comes from the type-specific congestion-cost coefficient rather than from type-specific conversion behaviour. We are explicit about this both here and in Section~\ref{sub:cb_decomposition}: the heterogeneous-type framework is what gives the model the parsimony to attach a corridor-specific cost coefficient to each type, but the calibrated model does not deliver a quantitatively significant class-specific conversion rate as an emergent output.

\paragraph{Policy implications for cross-border use.} The November 2025 Banco Central do Brasil package illustrates the policy use of the cost comparison. The package adds 150 basis points to the per-unit cost of using the stablecoin in the affected corridor. The net metric in the cross-border corridor moves from $-294$ basis points pre-package to $-144$ basis points post-package: the stablecoin remains cheaper by 144 basis points once the regulatory cost is internalised. For an authority calibrating the strength of such a package, the actionable mapping is direct. If the policy intent is to displace stablecoins, the package as designed is under-strength by roughly its own magnitude; an additional 144 basis points of friction, through tax wedges, reporting requirements, or restricted on-ramp access, would equalise the two costs. If the policy intent is to slow but not stop substitution, the package is already at the right magnitude. Section~\ref{sub:sens_brazil} reports the same comparison for shock magnitudes of 75, 150, 294 and 500 basis points, which gives a regulator a basis-point map between the chosen friction and the resulting equalisation point.

\paragraph{Implications for the retail corridor.} The retail corridor (US to Mexico) shows a negative net metric of 339 basis points at the calibrated parameters. The dominance is robust to the choice of risk-aversion weight $\gamma_k$ because the variance of $p^{sec}$ at the calibrated stress regime is small relative to the fee-cost difference. The stablecoin dominance in this corridor reflects two facts that the model is asked to take seriously: the incumbent fintech remittance fee averages 479 basis points in the RPW panel, and the stablecoin haircut conditional on redemption is bounded at roughly 240 basis points in the calibrated stress regime, weighted by the empirical regime frequency of 1.08\%. Even multiplying the stress frequency by an order of magnitude does not flip the sign.

\paragraph{The wholesale corridor as a counterfactual.} The wholesale corridor (Europe to Brazil) delivers a positive net metric of 613 basis points at the calibrated parameters, conditional on a wholesale congestion-cost coefficient of $c_q^{\text{wholesale}} = 1.00$. The coefficient is not derived from data; it is an input that translates the institutional requirement of atomic settlement into a basis-point cost. The wholesale result is therefore a counterfactual rather than a derived prediction: under the stated settlement-finality assumption, the stablecoin does not displace the incumbent rail. The same exercise with a retail-level congestion coefficient would flip the sign of the wholesale net metric, which we acknowledge by reporting the breakeven coefficient in the sensitivity appendix. The reading is that the cost comparison is informative for retail and cross-border use cases where settlement-finality requirements are mild, and that any extension to wholesale use cases needs the congestion-cost coefficient to be tied to data on TARGET2-level working-capital costs.

\paragraph{The boundary of the framework.} The calibrated parameters transport to DAI cleanly and do not transport to USDT, whose primary-market structure sits outside the model's open-arbitrage assumption. Section~\ref{sec:conclusion} discusses the boundary in detail; here we note only that policy use of the per-corridor cost comparison is reliable within the open-arbitrage segment of the stablecoin universe.

\paragraph{Open issues.} The calibration window is short, and the loss function uses six moments to identify a twelve-parameter vector. The eight structural parameters $(\sigma_k, \psi, \nu, \mu_q, \zeta)$ are weakly identified by price moments alone, with the regime means and precisions absorbing most of the empirical content. Richer on-chain data on per-type conversion timing would strengthen the identification of the type-specific signal noises $\sigma_k$ and provide independent moments for the calibration. The Brazilian framework is recent at the time of writing, and the post-package flow data is still thin; the comparison reported here therefore uses the pre-package window as the primary identification, with the post-package window as a robustness check. Extending the cost comparison to a multi-period setting with state evolution is described in the conclusion as the natural next step.
